



\chapter{CFEngine Standard Library}





\section{CFEngine Standard Library}

\begin{description}
\item[\textbf{COPBL}] Community Open Promise Body Library (aka CFEngine Standard Library)
\end{description}





CFEngine ships with a standard library of promise bodies and bundles
dealing with common aspects of system administration.

The CFEngine Standard Library is growing to include all common aspects
of system administration.

\begin{longtable}{|l|l|l|}
\hline
CFEngine version & Promise bodies & Promise bundles\\
\hline
\emph{3.1.5} & 88 & ?\\
\emph{3.2.1} & 99 & 19\\
\emph{3.3.5} & 114 & 29\\
\emph{3.3.8} & 113 & 26\\
\emph{3.4.4} & 124 & 32\\
\emph{3.10.3} & 189 & 127\\
\emph{3.12.6} & 193 & 134\\
\hline
\end{longtable}

\begin{codelisting}
\codecaption{420-060-COPBL-0480-File\_Exists\_And\_Is\_Mode\_612\_Without\_COPBL.cf}
%= lang:cfengine3, options: "linenos": false
\begin{code}
bundle agent main {

  files:

      "/tmp/testfile"
        comment => "Demonstrate setting file attributes",
        create  => "true",
        perms   => mog("612","aleksey","cfengine");

}

############################################################

body perms mog(mode,owner,group)
{
        owners => { "$(owner)", "john", "brian" };
        mode   => "$(mode)";
        groups => { "$(group)" };
}

\end{code}
\end{codelisting}
\begin{codelisting}
\codecaption{420-060-COPBL-0490-File\_exists\_and\_is\_mode\_6\_1\_2\_mog.cf}
%= lang:cfengine3, options: "linenos": false
\begin{code}
body file control
{
  inputs => { "$(sys.libdir)/stdlib.cf" };  # <-- read in the library
}

bundle agent main
{
  reports:
      "CFEngine StdLib is here: $(sys.libdir)/stdlib.cf" ; 

  files:
      "/tmp/testfile"
        create  => "true",
        perms   => mog("612","root","nobody");
}

\end{code}
\end{codelisting}
\begin{codelisting}
\codecaption{420-060-COPBL-0500-Context\_sensitive\_file\_editing\_Set\_robs\_password.cf}
%= lang:cfengine3, options: "linenos": false
\begin{code}
bundle agent main
{
  files:
      "/etc/shadow"
        handle => "context_sensitive_file_editing_demo",
        comment => "demonstrate context-sensitive file editing capability",
        edit_line => set_user_field("rob",
                                    "2",
                                    "$1$stIAaUZw$ptP75nVkz/EapeuvdWLNC0");
}

body file control { inputs => { "$(sys.libdir)/stdlib.cf" }; }
\end{code}
\end{codelisting}
\begin{codelisting}
\codecaption{420-060-COPBL-0510-Removing\_a\_file.cf}
%= lang:cfengine3, options: "linenos": false
\begin{code}
bundle agent main
{
  files:
      "/tmp/testfile.*"
        handle => "demo_removing_files",
        comment => "Demonstrate removing files using body delete tidy",
        delete => tidy;

      # shell equivalent:  rm -r /tmp/testfile*

}

body file control { inputs => { "$(sys.libdir)/stdlib.cf" }; }
\end{code}
\end{codelisting}





\begin{aside}
\label{aside:exercise_30}
\heading{Run the following command to generate some content:}

%= lang:bash
\begin{code}
date  > /tmp/date.txt
\end{code}

Write a CFEngine policy that will comment out (using \#)
all lines that start with a day of the week:

\begin{verbatim}
edit_line => comment_lines_matching("(Mon|Tue|Wed|Thur|Fri|Sat|Sun).*" ,
                                    "#")
\end{verbatim}

The \kode{edit\_line} bundle \kode{comment\_lines\_matching} is in the standard
library.

\end{aside}
\begin{codelisting}
\codecaption{420-060-COPBL-0530-Commenting\_out\_file\_contents.cf}
%= lang:cfengine3, options: "linenos": false
\begin{code}
bundle agent main
{
  files:
      "/etc/httpd/conf.d/maintenance.conf"
        handle    => "take_website_out_of_maintenance",
        comment   => "Disable maintenance-mode config block",
        edit_line => comment_out_everything;
}

bundle edit_line comment_out_everything
{
  replace_patterns:
      "^([^#].*)"
        replace_with => comment("# ");
}

body replace_with comment(c)
{
        replace_value => "$(c) $(match.1)";
        occurrences => "all";
}


\end{code}
\end{codelisting}
\begin{codelisting}
\codecaption{420-060-COPBL-0540-Z.cf}
%= lang:cfengine3, options: "linenos": false
\begin{code}
bundle agent main
{
  files:
      "/etc/httpd/conf.d/maintenance.conf"
        handle    => "put_website_into_maintenance",
        comment   => "Enable maintenance-mode config block",
        edit_line => uncomment_everything;
}

bundle edit_line uncomment_everything {
  replace_patterns:
      "^#(.*)"
        handle => "uncomment_everything_replace_pattern",
        comment => "If it starts with a hash mark, grab everything
                          after the hash mark, and uncomment it.",
        replace_with => uncomment;
}


body replace_with uncomment
{
        replace_value => "$(match.1)";
        occurrences => "all";
}


\end{code}
\end{codelisting}
\begin{codelisting}
\codecaption{420-060-COPBL-0550-Removing\_a\_file\_Remove\_centos\_httpd\_welcome\_page.cf}
%= lang:cfengine3, options: "linenos": false
\begin{code}
bundle agent main
{
  files:
      "/etc/httpd/conf.d/welcome.conf"
        handle => "nuke_welcome_conf",
        comment => "Let's keep a low profile and not advertise
                          what software we are running - remove the
                          Welcome page that says we are running Apache
                          on CentOS",
        delete => tidy;
}

body file control { inputs => { "$(sys.libdir)/stdlib.cf" }; }
\end{code}
\end{codelisting}
\begin{codelisting}
\codecaption{420-060-COPBL-0560-Remove\_httpd\_welcome\_page\_by\_commenting\_out\_welcome\_conf.cf}
%= lang:cfengine3, options: "linenos": false
\begin{code}
# welcome.conf is part of the Apache RPM
# to preserve package integrity, comment out this file's contents
# instead of deleting the file

bundle agent main
{
  files:
      "/etc/httpd/conf.d/welcome.conf"
        handle => "comment_out_welcome_dot_conf",
        comment => "Let's not ask for trouble by advertising
                          what software we are running",
        edit_line => comment_out_everything,
        classes => if_repaired("reload_httpd");

  commands:
    reload_httpd::
      "/etc/init.d/httpd"
        handle => "cmd_reload_httpd",
        comment => "Reload httpd configuration",
        args => "reload";
}

bundle edit_line comment_out_everything
{
  replace_patterns:
      "^([^#].*)"
        handle => "comment_out_everything_replace_patterns_promise",
        comment => "If it doesn't start with #, comment it out",
        replace_with => comment("#disabled-by-cfengine# ");
}

body file control { inputs => { "$(sys.libdir)/stdlib.cf" }; }
\end{code}
\end{codelisting}
\begin{codelisting}
\codecaption{420-060-COPBL-0590-classes\_if\_else.cf}
%= lang:cfengine3, options: "linenos": false
\begin{code}
bundle agent main
{
  files:
      "/dev/null/motd"
        handle => "touch_file",
        comment => "Demonstrate body classes if_else",
        create => "true",
        classes => if_else("file_exists","file_missing");

  reports:
    file_exists::
      "All OK"
        handle => "report_OK";

  reports:
    file_missing::
      "WARNING! Unable to create vital file!"
        handle => "report_WARN";
}

body file control { inputs => { "$(sys.libdir)/stdlib.cf" }; }
\end{code}
\end{codelisting}
\begin{codelisting}
\codecaption{420-060-COPBL-0600-classes\_persistent.cf}
%= lang:cfengine3, options: "linenos": false
\begin{code}
# "body classes state_repaired" is in StdLib

bundle agent main
{
  files:
    !file_fixed::
      "/tmp/file.txt"
        handle => "persistent_class_demo",
        comment => "Set a persistent class",
        create => "true",
        classes  => state_repaired("file_fixed");

  reports:
    file_fixed::
      "Persistent class set.  Run in verbose mode to see TTL"
        handle => "report_success",
        comment => "Report if our persistent class persistent_class
                          has been set as expected.";
}

body classes state_repaired(x)
{
        promise_repaired => { "$(x)" };
        persist_time => "10";
}
\end{code}
\end{codelisting}
\begin{codelisting}
\codecaption{420-060-COPBL-0610-comment\_lines\_matching.cf}
%= lang:cfengine3, options: "linenos": false
\begin{code}
bundle agent main
{
  files:
      "/tmp/scratch"
        handle => "selective_commenting",
        comment => "Remove specific lines",
        edit_line => comment_lines_matching("hello world", "#");
}


body file control { inputs => { "$(sys.libdir)/stdlib.cf" }; }
\end{code}
\end{codelisting}
\begin{codelisting}
\codecaption{420-060-COPBL-0620-contain\_silent.cf}
%= lang:cfengine3, options: "linenos": false
\begin{code}
bundle agent main
{
  commands:
      "/bin/date"
        handle => "run_date_cmd",
        comment => "Demonstrate 'body contain silent'",
        contain => silent;
}

body file control { inputs => { "$(sys.libdir)/stdlib.cf" }; }
\end{code}
\end{codelisting}
\begin{codelisting}
\codecaption{420-060-COPBL-0630-edit\_line\_insert\_lines.cf}
%= lang:cfengine3, options: "linenos": false
\begin{code}
bundle agent main
{
  files:
      "/etc/profile"
        handle => "edit_etc_profile",
        create => "true",
        edit_line => insert_lines("export ORGANIZATION=ACME");
}

body file control { inputs => { "$(sys.libdir)/stdlib.cf" }; }
\end{code}
\end{codelisting}
\begin{codelisting}
\codecaption{420-060-COPBL-0640-edit\_resolv\_dot\_conf\_using\_COPBL.cf}
%= lang:cfengine3, options: "linenos": false
\begin{code}
bundle agent main
{
  vars:
      "search_suffix"  string =>  "example.com example2.com";
      "nameservers"    slist  =>  { "8.8.4.4", "8.8.8.8" };

  files:
      "/tmp/resolv.conf"
        handle =>  "edit_resolv_conf",
        comment => "Setup up DNS resolver",
        edit_line =>  resolvconf("$(search_suffix)",
                                 "@($(this.bundle).nameservers)" );
}

body file control { inputs => { "$(sys.libdir)/stdlib.cf" }; }
\end{code}
\end{codelisting}
\begin{codelisting}
\codecaption{420-060-COPBL-0660-insert\_lines.cf}
%= lang:cfengine3, options: "linenos": false
\begin{code}
bundle agent main
{
  files:
      "/tmp/scratch"
        handle => "files_multi_line_insert",
        comment => "Insert multi-line content",
        create => "true",
        edit_line => insert_lines("
hello world
this is line 2
line 3 is great
line 4 is awesome
");
}


body file control { inputs => { "$(sys.libdir)/stdlib.cf" }; }
\end{code}
\end{codelisting}
\begin{codelisting}
\codecaption{420-060-COPBL-0680-set\_variable\_values.cf}
%= lang:cfengine3, options: "linenos": false
\begin{code}
bundle common g
{
  vars:

      "stuff[location]"     string => "New York";
      "stuff[time]"         string => "May-2027";
      "stuff[students]"     string => "11";
      "stuff[lab]"          string => "true";
}

bundle agent main {

  files:

      "/etc/example.conf"
        handle => "populate_config_file_from_array",
        comment => "Demonstrate 'bundle edit_line set_variable_values'",
        create => "true",
        edit_line => set_variable_values("g.stuff");
}


body file control { inputs => { "$(sys.libdir)/stdlib.cf" }; }
\end{code}
\end{codelisting}
\begin{codelisting}
\codecaption{420-060-StdLib-0470-Package\_add\_using\_StdLib.cf}
%= lang:cfengine3, options: "linenos": false
\begin{code}
bundle agent main
{
  packages:
      "php-mysql"
        policy => "present";
}

body common control {
# the following promise attributes support the 3.7 packages promises

  debian|redhat::
    package_inventory => { $(package_module_knowledge.platform_default) };
    package_module => $(package_module_knowledge.platform_default);
}

body file control { inputs => { "$(sys.libdir)/stdlib.cf" }; }

# Example run:
# root@localhost# cf-agent -f ~/p.cf -IC
# Warning: RPMDB altered outside of yum.
# ** Found 1 pre-existing rpmdb problem(s), 'yum check' output follows:
# 1:gdm-plugin-fingerprint-2.30.4-64.el6.x86_64 has missing requires of fprintd-pam
#     info: Successfully installed package 'php-mysql'
# root@localhost#
\end{code}
\end{codelisting}